\section{Execution of modified guest OS's}
One method of execution of modified guest OS’s is by
virtualisation at the operating system level. Using this
method, the physical server is virtualised at the operating
system level. In this case, the host operating system is a
modified kernel that permits the operation of several
isolated containers, sometimes referred to as virtualised
servers, virtual private servers, or when it comes to the
world of BSD, they are called jails. Every container is an
instance that uses the same host operating system kernel.
OpenVZ, Solaris Zones, and Linux-VServer are a few
examples that are using this method. The
implementations of this technology are widely utilised
and exhibit low overhead. The primary disadvantage of
para-virtualization is that it does not support multiple
kernels. Paravirtualization is the method that substitutes
the instructions of the actual machine's instruction set
architecture with a unique set of instructions known as
Hypercalls. Its low virtualisation overhead is a plus.
Denali \cite{citation_11}, Xen \cite{citation_5}, and Hyper-V \cite{citation_12} are a few
examples of para-virtualization solutions. In order to
allow guest operating systems to operate at higher levels,
the VMM (also known as the Hypervisor) in the x86
architecture operates directly above the actual hardware
(Ring 0). The main disadvantage of this strategy, despite
its support for multiple kernels, is that the guest OS’s
kernels must be altered in order to fully utilise the
Hypercalls. 